\documentclass[exercice,a4]{thedocument}%
\usepackage[french]{babel}%
\usepackage{pstricks}%
\usepackage{pst-plot}%
\usepackage{multido}%
\usepackage{tabularx}%
\usepackage{array}%
\usepackage[np]{numprint}%
\startdatakeys
    \source{22GENMATMEAG1}
    \BO{2026}
    \cycle{4}
    \competencesDuSocle{CA,RA}
    \objectifsApprentissage{C4C2ab,C4D3bf,C4B1ed}
    \motsCles{transformation,antécédent,fonction,médiane,agrandissement,aire}
\enddatakeys
\begin{document}%

\medskip
Cet exercice est un questionnaire à choix multiples (QCM). Aucune justification n'est demandée.

Pour chaque question, trois réponses (A, B et C) sont proposées. \textbf{Une seule réponse est exacte.}

Recopier sur la copie le numéro de la question et la réponse.

\medskip

{\footnotesize
\begin{tabularx}{\linewidth}{|m{8cm}|*{3}{>{\centering \arraybackslash}X|}}\hline
\centering{Questions} & Réponse A & Réponse B & Réponse C \\
\hline
\begin{minipage}{7.5cm}
\textbf{1.} On considère les deux figures suivantes.

Par quelle transformation la figure 2 est-elle l'image de la figure 1 ?

\psset{xunit=1.cm,yunit=1.cm,labelFontSize=\scriptstyle}
\begin{pspicture}(-1,-0.5)(5,3.5)
\pscustom[fillcolor=lightgray,fillstyle=solid]{\psline(0,2)(0,3)\psline(0,3)(1,3)\psline(1,3)(1,2)\psline(1,2)(2,1)\psline(2,1)(1,1)\psline(1,1)(0,2)}
\pscustom[fillstyle=solid,fillcolor=lightgray]{\psline(3,1)(3,2)\psline(3,2)(4,2)\psline(4,2)(4,1)\psline(4,1)(5,0)\psline(5,0)(4,0)\psline(4,0)(3,1)}
\multido{\n=0+1}{6}{\psline[linewidth=0.75pt,linecolor=gray](\n,-0.4)(\n,3.2)}
\multido{\n=0+1}{4}{\psline[linewidth=0.75pt,linecolor=gray](-0.2,\n)(5,\n)}
\uput[ur](1,3){B}\uput[ur](1,2){A}\uput[dr](1,1){E}
\uput[ur](4,2){B'}\uput[ur](4,1){A'}\uput[dr](4,0){E'}
\psline[linewidth=0.15pt,linecolor=lightgray]{->}(1,3)(4,2)
\psline[linewidth=0.15pt,linecolor=lightgray]{->}(1,2)(4,1)
\psline[linewidth=0.15pt,linecolor=lightgray]{->}(1,1)(4,0)
\uput[u](0.5,2){\huge\ding{192}} \uput[u](3.5,1){\huge\ding{193}}
\end{pspicture}
\end{minipage}
&une translation & une homothétie & une symétrie axiale\\\hline
\begin{minipage}{5cm}
\textbf{2.} On considère la représentation graphique de la fonction $g$ suivante :

\psset{unit=0.8cm,labelFontSize=\scriptstyle}
\begin{pspicture}(-3,-1)(4,5.2)
\multido{\n=-1+1}{6}{\psline[linewidth=0.75pt,linecolor=lightgray](\n,-1)(\n,5.1)}
\multido{\n=-1+1}{7}{\psline[linewidth=0.75pt,linecolor=lightgray](-1.1,\n)(4.1,\n) }
\psaxes[linewidth=0.95pt]{->}(0,0)(-1.1,-1.1)(4.1,5.1)
\psplot[linewidth=1.5pt,linecolor=blue,plotpoints=5000]{-0.50}{2.5}{x 2  mul }
\end{pspicture}

Quel est l'antécédent de 2 par la fonction $g$.
\end{minipage}
&2&1&4\\\hline

\begin{minipage}{7.5cm}

\smallskip
\textbf{3.} Soit $f$ la fonction définie par :
\[f(x)= 3x^2-7 \]
Quelle affirmation est correcte ?
\end{minipage}
&
 29 est l'image de 2 par la fonction $f$& $f(3) = 20$&$f$ est une fonction affine.\\
\hline

\begin{minipage}{7.5cm}
\smallskip
\textbf{4.} On a relevé les performances, en mètres, obtenues au lancer du poids
 par un groupe de 13 élèves d'une classe.
\np[m]{3,41} ; \np[m]{5.25} ; \np[m]{5.42} ; \np[m]{4.3} ; \np[m]{6.11} ;
\np[m]{4.28} ; \np[m]{5,15} ; \np[m]{3,7} ; \np[m]{6,07} ;
\np[m]{5,82} ; \np[m]{4,62} ; \np[m]{4,91} ; \np[m]{4,01}

Quelle est la médiane de cette série de valeurs ?
\end{minipage}
& 7&\np{4.91}&\np{5.15}\\
\hline
\begin{minipage}{7.5cm}

\smallskip
\textbf{5.} On considère la configuration suivante, dans laquelle les
triangles LAC et BUT sont semblables.

\psset{xunit=0.5cm,yunit=0.5cm,linewidth=0.5pt}
\begin{pspicture}(-1.5,-0.5)(7,6.75)
\psline(0,3)(2.49,2.84)\psline(0,3)(1.1,4.79)\psline(1.1,4.79)(2.49,2.84)
\psline(3.7,0)(6.99,5.37)\psline(3.7,0)(11.18,-0.48)\psline(11.18,-0.48)(6.99,5.37)
\uput[l](0,3){L}\uput[u](1.1,4.79){A}\uput[r](2.5,2.9){C}
\uput[l](3.7,0){B}\uput[u](6.99,5.37){U}\uput[r](11.18,-0.48){T}
\labelsep=1pt
\uput[l](0.5,3.7){\np[cm]{2.1}}\uput[dl](2,2.5){\np[cm]{2.5}}\uput[ur](2,3.5){\np[cm]{2.4}}
\uput[dr](3.8,2.7){\np[cm]{6.3}}\uput[dl](6.5,-0.45){\np[cm]{7.5}}\uput[ur](8.6,3.1){\np[cm]{7.2}}
\end{pspicture}

\medskip
Par quel nombre doit-on multiplier l'aire du triangle LAC pour obtenir l'aire du triangle BUT ?
\end{minipage}
&3&6&9\\\hline
\end{tabularx}
}
\bigskip

\showthesource
\end{document}
